\bookchapter{第二章：デジタルルームを整理する——ファイル管理}

\chapterintro{スマートフォンを長く使うと乱雑になる。良い習慣はファイル管理から始まる。部屋を定期的に整理する必要があるのと同じように、デジタル空間も手入れが必要だ。整然としたスマートフォンは、物を探しやすいだけでなく、動作もより滑らかになる。}

\booksection{一、3層ストレージ体系：あなたのデジタル倉庫}

\begin{enumerate}
  \item \textbf{本体ストレージ}：コア領域で、速度が速く、システムとアプリがここにある。スマートフォンで最も貴重な空間で、慎重に管理すべき。
  \item \textbf{SDカード（サポートの場合）}：拡張クローゼットで、速度は遅いが、映画、写真などの大容量ファイルに適する。注意：現在多くのフラッグシップ機はSDカードをサポートしない。
  \item \textbf{クラウドストレージ}：金庫＋同期ドライブ。紛失防止と複数デバイス間同期が可能。iCloud、Baidu Netdisk、OneDriveなどがよく使われる。
\end{enumerate}

\booksection{二、「ファイルマネージャー」を知る}

Androidには通常 [ファイル] または [ファイル管理] アプリがあり、iOSは標準の [ファイル] アプリだ。これらはスマートフォンのすべてのファイルを表示する「エクスプローラー」だ。

\booksubsection{一般的なフォルダの意味：}

\begin{itemize}
  \item \textbf{DCIM}：カメラで撮影した写真と動画。
  \item \textbf{Pictures}：その他の画像、スクリーンショット、保存した画像など。
  \item \textbf{Download}：ブラウザでダウンロードしたファイル。
  \item \textbf{Music}：音楽ファイル。
  \item \textbf{Movies}：動画ファイル。
  \item \textbf{Documents}：ドキュメントファイル。
  \item \textbf{Android}：Androidシステムフォルダで、中に各アプリのデータがある。通常は手動削除しないこと。
\end{itemize}

\booksection{三、ファイルシステム構造：より深く理解する}

AndroidはLinuxベースで、ファイルシステムには一定の階層構造がある。一般ユーザーが深く理解する必要はないが、基本的な概念を知っておくと役立つ。

\booksubsection{ストレージパーティション：}

\begin{itemize}
  \item \textbf{/system}：システムパーティションで、オペレーティングシステムファイルを格納。一般ユーザーはアクセス権限がなく、root権限が必要。
  \item \textbf{/data}：データパーティションで、ユーザーデータとアプリデータを格納。日常使用する空間。
  \item \textbf{/sdcard または /storage/emulated/0}：ユーザーがアクセス可能なストレージ空間で、写真やダウンロードファイルはここにある。
\end{itemize}

\booksubsection{アプリデータストレージ：}

\begin{itemize}
  \item \textbf{内部ストレージ}：アプリのプライベート空間で、他のアプリはアクセス不可。アプリアンインストール時に削除される。
  \item \textbf{外部ストレージ}：アプリがユーザーアクセス可能領域に保存するファイル、例えば写真、ダウンロードコンテンツ。アプリアンインストール時に必ずしも削除されない。
\end{itemize}

\booksection{四、データバックアップと復元：データに保険をかける}

データはプライスレス、バックアップは必須の習慣だ。

\booksubsection{バックアップ戦略：}

\begin{enumerate}
  \item \textbf{クラウドバックアップ}：自動で便利だが、ネットワークとクラウド容量が必要。iCloud、Huawei Cloud、Xiaomi Cloudなどがサポート。
  \item \textbf{ローカルバックアップ}：データケーブルでパソコンに接続し、公式ツールでバックアップ。速く容量が大きいが、手動操作が必要。
  \item \textbf{二重バックアップ}：クラウドバックアップ＋ローカルバックアップで最も安全。
\end{enumerate}

\booksubsection{何をバックアップする？}

\begin{itemize}
  \item \textbf{写真と動画}：最も容量を占有し、最も貴重。
  \item \textbf{連絡先}：クラウドアカウントに保存するのが最も便利。
  \item \textbf{WeChatチャット履歴}：WeChat内蔵のバックアップ機能を使用。
  \item \textbf{重要ドキュメント}：クラウドドライブやパソコンに保存。
\end{itemize}

\booksection{五、ファイルセキュリティ：プライバシーを守る}

\begin{itemize}
  \item \textbf{センシティブファイルの暗号化}：多くのスマートフォンには「セキュアフォルダ」や「プライベート空間」機能があり、センシティブファイルを入れ、パスワードや指紋でアクセスできる。
  \item \textbf{削除$\neq$抹消}：通常の削除は空間が上書き可能というマークをつけるだけで、データは残っている。完全に削除するには、専用の安全消去ツールを使用するか、複数回上書き書き込みが必要。
  \item \textbf{クラウドセキュリティ}：信頼できるクラウドサービスプロバイダーを選び、二要素認証を有効化し、定期的にログインデバイスを確認。
\end{itemize}

\booksection{六、容量が足りない？大掃除をしよう}

\begin{itemize}
  \item \textbf{キャッシュクリーン}：アプリの一時記憶で、クリーンすると空き容量が増える。設定 > ストレージでクリーン可能。
  \item \textbf{大容量ファイルの検索}：ファイルマネージャーでサイズ順にソートし、容量を占有する「犯人」を見つける。
  \item \textbf{残存物の削除}：アプリアンインストール後、ゴミファイルが残ることがあるので手動で削除。
  \item \textbf{アルバムのクリーン}：重複写真、ブレた写真、不要なスクリーンショットを削除。
  \item \textbf{WeChatのクリーン}：WeChatは容量の大口で、定期的なチャット履歴とキャッシュのクリーンが必須。
\end{itemize}

\booksection{七、一般的なファイル形式の全解析}

スマートフォン内の各種ファイル拡張子は複雑だが、理解すればファイル管理がより上手になる：

\begin{center}
\begin{tabularx}{\linewidth}{|X|X|X|X|}
\hline
\textbf{形式} & \textbf{タイプ} & \textbf{特性} & \textbf{推奨される開き方} \tabularnewline
\hline
JPG/JPEG & 画像 & 非可逆圧縮、サイズ小、互換性最高。スマートフォン撮影のデフォルト形式。 & アルバム/任意の画像ビューアー \tabularnewline
PNG & 画像 & 可逆圧縮、透明背景サポート。スクリーンショット、アイコンに常用。JPGよりサイズ大。 & アルバム/任意の画像ビューアー \tabularnewline
MP3 & オーディオ & 非可逆圧縮、サイズ小、互換性極めて高い。最も汎用的なオーディオ形式。 & 任意の音楽プレーヤー \tabularnewline
FLAC & オーディオ & 可逆圧縮、音質はCDと同等、サイズはMP3の3-5倍。オーディオマニアの第一選択。 & 音楽プレーヤー/プロフェッショナル再生アプリ \tabularnewline
WAV & オーディオ & 無圧縮の生オーディオ、音質最高だがサイズ巨大。プロ録音に常用。 & 音楽プレーヤー \tabularnewline
M4A & オーディオ & AACエンコード、MP3の同ビットレートより音質が良い。Appleエコシステムで常用。 & Apple Music/音楽プレーヤー \tabularnewline
MP4 & 動画 & 最も汎用的な動画コンテナ形式、互換性極めて高い。スマートフォン録画のデフォルト形式。 & アルバム/任意の動画プレーヤー \tabularnewline
MKV & 動画 & オープンソースコンテナ形式、マルチオーディオトラック、マルチ字幕サポート。映画リソースに常用。 & VLC/nPlayer \tabularnewline
DOC/DOCX & ドキュメント & Wordドキュメント。DOCXは圧縮形式でサイズがより小さい。DOCは旧形式。 & WPS/Office \tabularnewline
PPT/PPTX & プレゼン & PowerPointプレゼン資料。PPTXは圧縮形式。 & WPS/Office \tabularnewline
XLS/XLSX & スプレッドシート & Excelスプレッドシート。XLSXは圧縮形式。 & WPS/Office \tabularnewline
PDF & ドキュメント & クロスプラットフォーム形式、レイアウト固定で編集不可。契約書、論文、電子書籍に常用。 & Adobe Reader/WPS \tabularnewline
ZIP & 圧縮 & 最も汎用的な圧縮形式、クロスプラットフォーム互換性良好。 & システム標準/解凍アプリ \tabularnewline
RAR & 圧縮 & 圧縮率はZIPより高いが、専用ソフトでの解凍が必要。 & 解凍アプリ \tabularnewline
ISO & イメージ & 光ディスクイメージファイル、光ディスクの内容を完全にコピー。システムインストーラに常用。 & 解凍アプリ/仮想光ドライブ \tabularnewline
HTML & ウェブページ & ウェブページファイル、ブラウザで開く。保存したウェブページに常用。 & ブラウザ \tabularnewline
\hline
\end{tabularx}
\end{center}

\begin{tcbpractice}{練習：「ダウンロード」フォルダを整理}
ファイルマネージャーを開き > 「ダウンロード」フォルダに入る。中に何があるか？不要なものを削除し、有用なものを分類（例えば「仕事資料」「スタンプ」フォルダを作成して入れる）。ついでにストレージ容量を確認し、どのタイプのファイルが最も多く占有しているか把握する。

\end{tcbpractice}

\booksection{八、画像とファイルシステム管理：資産を整然と}

前節で「フォルダを知る」ことを学んだが、知るだけでは不十分。本節では\textbf{長期的に実行可能な管理規範を主体的に構築する方法}を説明し、ファイルがいつでも見つかり、失われず、覗き見されないようにする。

\booksubsection{ファイル分類と命名規範}

混乱の根源はルールがないこと。ルールの構築は実は簡単で、重要なのは\textbf{継続}。

\begin{itemize}
  \item \textbf{「時間＋イベント」命名法}：フォルダ名にタイムスタンプとイベント記述を含め、後日の回想と検索に便利。例えば\texttt{202606\_旅行チケット}、\texttt{2025年年末まとめ}、\texttt{2024子供の誕生日}。これでタイムラインが見え、内容も一目で分かる。
  \item \textbf{3階層分類体系}：
  \begin{itemize}
    \item \textbf{第1層}：ファイルタイプ別（ドキュメント / 画像 / 動画 / オーディオ）
    \item \textbf{第2層}：使用シーン別（仕事 / 生活 / 学習）
    \item \textbf{第3層}：具体的なイベントやプロジェクト別（例：「2026Q2\_プロジェクトA」）
  \end{itemize}
  \item \textbf{月次クリーン計画}：毎月固定日付（給料日や月初1日を推奨）に、「ダウンロード」ディレクトリを一括整理：\textbf{有用なものをアーカイブ、不要なものを削除、貴重なものを転送保存}。年12回の小整理は、1年分の大掃除よりはるかに楽。
\end{itemize}

\begin{tcbpractice}{今月のタスクリスト}
暇な週末を選び、以下の順序で操作：

\begin{enumerate}
  \item ファイルマネージャーで\texttt{2026Q2\_プロジェクトA}形式の今月の仕事フォルダを作成；
  \item 「ダウンロード」ディレクトリ内の今週のファイルをすべて確認し、有用なものを残し、アーカイブし、不要なものを削除；
  \item 今年の家族写真に\texttt{2025家族}、\texttt{2026家族}という年次フォルダを作成。
\end{enumerate}

\end{tcbpractice}

\booksubsection{冗長ファイル処理戦略}

スマートフォンを長く使うと「太る」が、多くは冗長ファイルの仕業だ。

\begin{itemize}
  \item \textbf{72時間自動クリーン}：WeChat、QQなどソーシャルアプリが受信する一時ファイル（スタンプ、チャット画像、ショート動画）は、ほとんど一度見たら不要。[WeChat] > [設定] > [一般] > [ストレージ] で「自動クリーン」を有効化し、期限切れキャッシュを72時間に設定。スマートフォンが瞬時にさっぱりし、重要な記録も失われない。
  \item \textbf{ストレージ容量の指向性クリーン}：Android [設定] > [ストレージ] > [クリーン加速]、iOS [設定] > [一般] > [iPhoneストレージ]で、システムが「キャッシュ / 一時ファイル / アンインストール可能アプリ / 大容量ファイル」別にリスト化。重点クリーン：(1) キャッシュファイル；(2) 期限切れインストーラ（アプリアップグレード後に残る旧版インストーラ）；(3) 重複ファイル（同名同サイズの写真）。
  \item \textbf{自動分類ルール}：大部分のブランドのアルバムアプリはスマート分類をサポート：人物、風景、ドキュメント、スクリーンショット。有効化すると、新しく撮った集合写真は自動的に「人物 > 張三」に分類され、スキャンした契約書は自動的に「ドキュメント」に分類される。\textbf{月次整理と組み合わせることで、「物が見つからない」確率を大幅に低減}。
\end{itemize}

\booksubsection{画像安全処理規範}

画像はスマートフォンで最も「センシティブ」なデータタイプだ——何気なく撮った1枚であなたの位置、デバイス、時間が漏れる。私たちは\textbf{「二重保護」ワークフローの構築}：

\begin{tcbwarn}{証明写真/センシティブスクリーンショット送信の標準フロー}
\begin{enumerate}
  \item \textbf{可視情報のマスキング}：システム標準またはサードパーティツール（Meitu XiuXiu、Snapseedなど）で\textbf{身分証番号の中間数桁、ナンバープレート、部屋番号、氏名、金額}などのキーフィールドを塗りつぶす；
  \item \textbf{メタデータ剥離}：送信前にアルバムで「共有」または「エクスポート」をタップすると、システムが位置情報を保持するか尋ねるので\textbf{「保持しない」}を選択。これで撮影デバイスの型番、GPS座標、撮影時間が画像と一緒に送信されない；
  \item \textbf{一時透かし}（オプション）：センシティブスクリーンショットに「XX手続き専用」のテキスト透かしを追加し、悪用を防止。
\end{enumerate}

\end{tcbwarn}

メタデータ（EXIF）削除の一般的な方法（基本操作、詳細な保護は第六章「写真のセンシティブ情報処理」を参照）：

\begin{itemize}
  \item \textbf{システム標準}：iOS [設定] > [プライバシーとセキュリティ] > [カメラ]で「正確な位置」をオフにできる；共有時に「現在の写真のみ」を選択すると\textbf{大部分のEXIFを自動剥離}。
  \item \textbf{Android}：大部分のブランドのアルバムアプリは「共有」時に「位置情報を削除するか」を確認し、チェックを入れるだけ。
  \item \textbf{プロフェッショナルツール}：ExifTool（PC版）、Metapho（iOS）、Photo Metadata Remover（Android）ですべてのメタデータを完全に削除可能。
\end{itemize}

長期保存\textbf{が必要なセンシティブ画像（身分証、銀行カード、戸籍謄本のスキャン、パスポート）は、}暗号化アルバム\textbf{に保存することを推奨：}：

\begin{itemize}
  \item Huawei/Xiaomi：[アルバム] > 写真を長押し > 「セキュアボックスに移動」；
  \item サードパーティ暗号化アプリ：
  \item 360セキュアアルバム\textbf{、}、\textbf{、}、\textbf{など、単独パスワード保護で隠蔽性がより高い。}第三章：スマートフォンに「生命線」をつなぐ——ネットワーク
\end{itemize}
